% ============================================================================
%  CROSS-EVAL METROLOGICAL SUMMARY  --  TEMPLATE ENGINE
%  ----------------------------------------------------------------------------
%  One model, many benchmarks. Compares metrological quality across evals.
%  Do NOT edit this file for a normal summary -- put values in
%  cross-eval-config.tex.
%
%  Build:   pdflatex cross-eval.tex   (run twice; or: make)
%
%  Inputs (all optional except the config):
%    cross-eval-config.tex   required  -- the parameter file
%    evals.csv               optional  -- per-eval metrics table
% ============================================================================
\documentclass[10pt]{article}

\usepackage[a4paper,margin=20mm,headheight=14pt]{geometry}
\usepackage[T1]{fontenc}
\IfFileExists{lmodern.sty}{%
  \usepackage{lmodern}%
  \IfFileExists{microtype.sty}{\usepackage{microtype}}{}%
}{%
  % No scalable Latin Modern: skip microtype to avoid bitmap-font errors.
}
\usepackage{xcolor}
\usepackage{array}
\usepackage{tabularx}
\usepackage{booktabs}
\usepackage{longtable}
\usepackage{enumitem}
\usepackage{titlesec}
\usepackage{fancyhdr}
\usepackage{tikz}
\usepackage{csvsimple}
\usepackage[hidelinks]{hyperref}

% ---------------------------------------------------------------- palette ----
\definecolor{rcink}{HTML}{1A1A1A}
\definecolor{rcaccent}{HTML}{8B4513}
\definecolor{rcrule}{HTML}{C9C2B6}
\definecolor{rcsoft}{HTML}{F4F1EA}
\definecolor{rcmute}{HTML}{6B6B6B}
\definecolor{rctier-strong}{HTML}{2E7D32}
\definecolor{rctier-mid}{HTML}{E65100}
\definecolor{rctier-fragile}{HTML}{C62828}
\color{rcink}

% ============================================================================
%  PARAMETER STORE  (identical engine to the single run card)
% ============================================================================
\makeatletter
\newcommand{\rcset}[2]{\expandafter\def\csname rc@@\detokenize{#1}\endcsname{#2}}
\newcommand{\rc}[1]{%
  \ifcsname rc@@\detokenize{#1}\endcsname
    \expandafter\let\expandafter\rc@cur\csname rc@@\detokenize{#1}\endcsname
    \ifx\rc@cur\@empty\textcolor{rcmute}{\textemdash}%
    \else\rc@cur\fi
  \else\textcolor{rcmute}{\textemdash}\fi}

\newcommand{\rchist}[1]{%
  \ifcsname rc@@\detokenize{#1}\endcsname
    \edef\rc@list{\csname rc@@\detokenize{#1}\endcsname}%
  \else\let\rc@list\@empty\fi
  \ifx\rc@list\@empty\textcolor{rcmute}{\textit{(no histogram data)}}\else
  \begin{tikzpicture}[baseline=-2pt]
    \pgfmathsetmacro{\rcmax}{0}%
    \foreach \v in \rc@list{%
      \pgfmathsetmacro{\rcmax}{max(\rcmax,\v)}\xdef\rcmax{\rcmax}}%
    \def\rcx{0}%
    \foreach \v in \rc@list{%
      \pgfmathsetmacro{\rcbh}{(\rcmax>0)*(\v/\rcmax)*1.25 + 0.012}%
      \fill[rcaccent] (\rcx,0) rectangle ++(0.30,\rcbh);%
      \pgfmathsetmacro{\rcx}{\rcx+0.36}\xdef\rcx{\rcx}}%
    \draw[rcrule,line width=0.4pt] (-0.02,0) -- (\rcx,0);%
  \end{tikzpicture}\fi}

% --------------------------------------------------------- section styling ----
\titleformat{\section}{\normalfont\large\bfseries\color{rcink}}{}{0pt}%
  {\rule[-3pt]{0pt}{14pt}}[{\color{rcrule}\titlerule[1pt]}]
\titlespacing*{\section}{0pt}{16pt}{8pt}
\setlist{nosep,leftmargin=1.4em}
\renewcommand{\arraystretch}{1.28}

\newenvironment{rcfacts}{\par\smallskip}{\par}
\newcommand{\fact}[2]{%
  \noindent
  \begin{minipage}[t]{0.25\linewidth}\raggedright\strut
    \bfseries\color{rcmute}#1\end{minipage}%
  \hfill
  \begin{minipage}[t]{0.71\linewidth}\raggedright\strut #2\end{minipage}%
  \par\smallskip}

\newcommand{\stat}[2]{%
  \begin{tabular}[t]{@{}l@{}}%
  {\footnotesize\color{rcmute}#1}\\[1pt]{\large #2}%
  \end{tabular}}

% ---------------------------------------------------- header / footer ----
\pagestyle{fancy}\fancyhf{}
\renewcommand{\headrulewidth}{0.4pt}
\renewcommand{\footrulewidth}{0pt}
\lhead{\footnotesize\color{rcmute}CROSS-EVAL METROLOGICAL SUMMARY}
\rhead{\footnotesize\color{rcmute}\texttt{\rc{summary.id}}}
\cfoot{\footnotesize\color{rcmute}\thepage}

% --------------------------------------------- disclaimer defaults ----
\@ifundefined{rc@@\detokenize{disc.singlemodel}}{\rcset{disc.singlemodel}{%
  \textbf{Single model.} All benchmarks were evaluated with one model and
  one serving configuration. Metrological quality rankings may differ for
  other models --- item sensitivity is model-dependent.}}{}
\@ifundefined{rc@@\detokenize{disc.perturbation}}{\rcset{disc.perturbation}{%
  \textbf{Perturbation is not validity.} Stability under cosmetic
  perturbation measures instrument reliability, not construct validity.
  A stable benchmark is not necessarily measuring anything meaningful;
  an unstable benchmark is definitely not.}}{}
\@ifundefined{rc@@\detokenize{disc.stopping}}{\rcset{disc.stopping}{%
  \textbf{Retrospective only.} Sequential-stopping figures are
  counterfactual replays over permutations of already-collected data.
  They are NOT a prospective error guarantee for a future early-stopped
  run.}}{}
\@ifundefined{rc@@\detokenize{disc.comparison}}{\rcset{disc.comparison}{%
  \textbf{Cross-eval comparison caveats.} Benchmarks differ in item count,
  domain, response format, and difficulty. Comparing metrological quality
  across benchmarks is informative but not controlled --- a larger item
  pool naturally admits more stable aggregate estimates.}}{}
\makeatother

% ============================================================================
\input{cross-eval-config.tex}

\begin{document}

% ---------------------------------------------------------------- TITLE ----
\begin{center}
  {\Huge\bfseries Cross-Eval Metrological Summary}\\[4pt]
  {\color{rcmute}\small One Model, Many Benchmarks --- Instrument Quality Comparison}
\end{center}
\vspace{4pt}
{\color{rcrule}\hrule height 1pt}
\vspace{6pt}
\noindent
\stat{SUMMARY ID}{\texttt{\rc{summary.id}}}\hfill
\stat{DATE ISSUED}{\rc{date.issued}}\hfill
\stat{N EVALS}{\rc{n.evals}}\hfill
\stat{MODEL}{\rc{model.short}}
\vspace{4pt}
{\color{rcrule}\hrule height 1pt}

% --------------------------------------------------------------- HEADER ----
\section{Header}
\begin{rcfacts}
\fact{Summary ID}{\texttt{\rc{summary.id}}}
\fact{Date issued}{\rc{date.issued}}
\fact{Model under test}{\rc{model.name}\\[2pt]
  \textcolor{rcmute}{revision:} \rc{model.revision} \quad
  \textcolor{rcmute}{quant:} \rc{model.quant}\\[2pt]
  \textcolor{rcmute}{serving:} \rc{model.serving}}
\fact{N evals}{\rc{n.evals}}
\fact{Eval selection}{\rc{eval.selection}}
\fact{Evaluator}{\rc{evaluator.org} \,\textbar\, \rc{evaluator.tool}}
\end{rcfacts}

% -------------------------------------------------- PERTURBATION DESIGN ----
\section{Perturbation Design \textnormal{\small(shared across evals)}}
\begin{rcfacts}
\fact{Design}{\rc{perturb.design}}
\fact{Inference config}{\rc{inference.config}}
\fact{N observations}{\rc{n.total.obs} \textnormal{(total item$\times$variant
  observations across all evals)}}
\end{rcfacts}

% ---- Per-eval accuracy + CI table -----------------------------------------
\section{Cross-Eval Accuracy}
\IfFileExists{evals.csv}{%
  \par\smallskip
  \csvreader[
    tabular={@{}l r r r r r r r@{}},
    respect all,
    table head={\toprule
      \bfseries Eval &
      \bfseries Acc. &
      \bfseries CI lo &
      \bfseries CI hi &
      \bfseries Spread &
      \bfseries Sens.\% &
      \bfseries Items &
      \bfseries Variants
      \\\midrule},
    table foot={\bottomrule},
    late after line=\\,
  ]{evals.csv}{}{\csvcoli & \csvcolii & \csvcoliii & \csvcoliv &
    \csvcolv & \csvcolvi & \csvcolvii & \csvcolviii}%
  \par\smallskip
  \begin{rcfacts}
  \fact{Accuracy range}{\rc{acc.best.name} (\rc{acc.best.val}) to
    \rc{acc.worst.name} (\rc{acc.worst.val})}
  \fact{Spread range}{\rc{spread.best.name} (\rc{spread.best.val}) to
    \rc{spread.worst.name} (\rc{spread.worst.val})}
  \end{rcfacts}
}{%
  \par\smallskip\textcolor{rcmute}{\itshape No \texttt{evals.csv} supplied.}\par}

% --------------------------------------------- METROLOGICAL QUALITY ----
\section{Metrological Quality}
\begin{rcfacts}
\fact{What this measures}{The fraction of items whose score flips under
  cosmetic perturbation (option reordering and temperature variation).
  Lower is better: a metrologically sound instrument gives the same
  answer regardless of incidental presentation choices.}
\end{rcfacts}

\par\medskip\noindent\textbf{\color{rcmute}Per-eval stability histograms}
{\small\color{rcmute}\ (fraction correct across variants, 10 bins 0--1)}
\par\smallskip
\begin{rcfacts}
\fact{\rc{eval.1.name}}{\rchist{eval.1.hist}\\[2pt]
  stable \rc{eval.1.stable} \quad
  sensitive \rc{eval.1.sensitive} (\rc{eval.1.pct}\%)}
\fact{\rc{eval.2.name}}{\rchist{eval.2.hist}\\[2pt]
  stable \rc{eval.2.stable} \quad
  sensitive \rc{eval.2.sensitive} (\rc{eval.2.pct}\%)}
\fact{\rc{eval.3.name}}{\rchist{eval.3.hist}\\[2pt]
  stable \rc{eval.3.stable} \quad
  sensitive \rc{eval.3.sensitive} (\rc{eval.3.pct}\%)}
\fact{\rc{eval.4.name}}{\rchist{eval.4.hist}\\[2pt]
  stable \rc{eval.4.stable} \quad
  sensitive \rc{eval.4.sensitive} (\rc{eval.4.pct}\%)}
\fact{\rc{eval.5.name}}{\rchist{eval.5.hist}\\[2pt]
  stable \rc{eval.5.stable} \quad
  sensitive \rc{eval.5.sensitive} (\rc{eval.5.pct}\%)}
\fact{\rc{eval.6.name}}{\rchist{eval.6.hist}\\[2pt]
  stable \rc{eval.6.stable} \quad
  sensitive \rc{eval.6.sensitive} (\rc{eval.6.pct}\%)}
\end{rcfacts}

% ------------------------------------------------ TIER CLASSIFICATION ----
\section{Tier Classification}
\begin{rcfacts}
\fact{Classification method}{\rc{tier.method}}
\fact{\textcolor{rctier-strong}{Strong}}{\rc{tier.strong}}
\fact{\textcolor{rctier-mid}{Middling}}{\rc{tier.middling}}
\fact{\textcolor{rctier-fragile}{Fragile}}{\rc{tier.fragile}}
\fact{Interpretation}{\rc{tier.interpretation}}
\end{rcfacts}

% ------------------------------------ SEQUENTIAL STOPPING ----
\section{Sequential Stopping \textnormal{\small(retrospective, per-eval)}}
\begin{rcfacts}
\fact{Method}{\rc{stop.method}}
\fact{Per-eval medians}{\rc{stop.summary}}
\fact{Interpretation}{\rc{stop.interpretation}}
\end{rcfacts}

% ----------------------------------------------------------- DISCLAIMERS ----
\section{Disclaimers}
\begin{itemize}[leftmargin=1.2em,itemsep=4pt]
  \item \rc{disc.singlemodel}
  \item \rc{disc.perturbation}
  \item \rc{disc.stopping}
  \item \rc{disc.comparison}
\end{itemize}

% ------------------------------------------------------ RAW DATA REFERENCE ----
\section{Raw Data Reference}
\begin{rcfacts}
\fact{Companion bundle}{\rc{companion.file}}
\fact{Hash}{\rc{companion.hash}}
\fact{Contents}{\rc{companion.contents}}
\end{rcfacts}

\end{document}
